% !TEX encoding = UTF-8 Unicode
% !TeX program = xelatex
% !TEX TS-program = XeLaTeX

\documentclass[12pt, a4paper]{article}
%\documentclass[12pt, a5paper]{article}
\usepackage{mathtools}
\usepackage{algorithm}
\usepackage[noend]{algpseudocode}
\usepackage{unicode-math}
\usepackage{comment}
\usepackage{wrapfig}
\usepackage{xfrac}
\usepackage{tikz}
\usetikzlibrary{arrows, decorations.markings}

\usepackage[pass]{geometry}
%\geometry{legalpaper, landscape, margin=1in}
\usepackage{fullpage}
\usepackage{setspace}
\usepackage[localise]{xepersian}
\setstretch{1.4}
\settextfont[Scale=1]{XB Niloofar}
\setlatintextfont[Scale=1]{XB Niloofar}%{Junicode} %Junicode
\setdigitfont{XB Niloofar}%{XB Zar}

\setlength{\parindent}{0ex}

\definecolor{answercolor}{HTML}{0080ff}
\newcommand{\vsp}{\vspace{34pt}}
%\newcommand{\پاسخ}[2]{#1}
%\thispagestyle{empty}
\newcommand{\پاسخ}[2]{\color{answercolor}\textbf{پاسخ: }#2\color{black}\vspace{-0.3cm}}

\begin{document}
		\tikzset{
		treenode/.style = {align=center, inner sep=0pt, text centered,
			font=\sffamily},
		arn_n/.style = {treenode, circle, white, font=\sffamily\bfseries, draw=black,
			fill=black, text width=1.5em},% arbre rouge noir, noeud noir
				normal/.style = {circle, white, font=\sffamily\bfseries, draw=black,
			fill=black, text width=1.5em},%
		arn_r/.style = {treenode, circle, red, draw=red, 
			text width=1.5em, very thick},% arbre rouge noir, noeud rouge
		arn_s/.style = {treenode, circle, black, draw=black, 
			text width=1.5em, very thick},% arbre rouge noir, noeud rouge
		arn_x/.style = {treenode, rectangle, draw=black,
			minimum width=0.5em, minimum height=0.5em, very thick}% arbre rouge noir, nil
	}
	\begin{tikzpicture}[remember picture,overlay]	\node[anchor=north,yshift=-1.5cm] at (current page.north) {\rl{به نام خدا}};
	\end{tikzpicture}%

\vspace{-1cm}
{\centering
	\textbf{کوئیز ششم درس ساختمان داده‌ها و الگوریتم‌ها}\\
مدرس: مهدی صفرنژاد       - ۴ خرداد ۹۸\\
\raggedleft\small
%نام و نام خانوادگی: \\
}
\vspace{-0.5cm}
\پرش‌متوسط\پرخط‌ا\\
در سؤالاتی که توضیح خواسته شده، یک پاسخ کوتاه (حدوداً بین یک تا دو خط) مدنظر است. در سؤالاتی که خواستۀ مسئله توضیحی نیست، جواب نهایی کفایت می‌کند و نوشتن مراحل رسیدن به جواب نیاز نیست.\\


\پاراگراف{سؤال ۱}
میخواهیم روی یک رشته $n$ بیتی از ۰ و ۱ الگوریتم زیر را اجرا کنیم. با توجه به این الگوریتم به سوالات زیر پاسخ دهید.
\begin{latin}
	\begin{algorithm}
		\begin{algorithmic}
			\While{$true$} 
			\State $cnt := 0$
			\For{$ i = 1 \to n-2$}
			\If{$s[i] = 1 \; and \; s[i+1] = 0$}
			\State $s[i] := 1 - s[i]$
			\State $s[i+1] := 1 - s[i + 1]$
			\State $cnt := cnt + 1$
			\EndIf
			\EndFor
			\If{$cnt = 0$}
			\State $break$
			\EndIf
			\EndWhile  \label{roy's loop}
		\end{algorithmic}
	\end{algorithm}
\end{latin}


\begin{itemize}
	\item 
	ثابت کنید این الگوریتم به پایان می‌رسد. [۲۰ نمره]
	\item
	پیچیدگی زمانی الگوریتم را در بهترین و بدترین حالت حساب کرده و برای هر کدام یک مثال بزنید.
	[۳۸ نمره]
\end{itemize}

\پاسخ{\vsp\vsp}{
	\begin{itemize}
		\item 
		اگر کل عدد را به فرم یک عدد باینری بزرگ ببینیم، کد بالا هر سری می‌بیند آیا ۱ و ۰ مجاوری وجود دارد یا نه و در صورت وجود آن‌ها را با هم جابه‌جا می‌کند. با انجام هر کدام از این تغییرات عدد باینری بزرگ کم می‌شود و می‌دانیم این عدد حداقل صفر می‌تواند باشد بنابراین فرآیند الگوریتم بالا پایان‌پذیر است.
		
		\item
		در بهترین حالت پیچیدگی الگوریتم
		$O(n)$
		است یعنی در این حالت هیچ جفتی برای تغییر نداریم. برای مثال فرض کنید تمام صفر ها در نیمه چپ و تمام ۱ ها در نیمه راست باشند.
		
		اما برای بدترین حالت نشان می‌دهیم تعداد جابه‌جایی‌ها برابر تعداد جفت‌ اندیس‌های 
		$i,j$
		است که 
		$i < j$
		و 
		$s[i] = 1$
		اما
		$s[j] = 0$
		توجه کنید در هر سری جابه‌جایی این تعداد جفت‌ها دقیقا یک واحد کم می‌شود و زمانی که به صفر برسد دیگر جابه‌جایی‌ای انجام نخواهد شد. بنابراین اگر تعداد این جفت‌ها بیشینه شود زمان اجرای الگوریتم نیز بیشینه خواهد شد. در حالتی بیشینه می‌شود که در نیمه سمت چپ تماما ۱ قرار بگیرد و در نیمه سمت راست تماما صفر که
		$\big(\frac{n}{2}\big)^2$
		تا از این جفت‌ها را می‌دهد و در بدترین حالت
		$\Theta(n^2)$
		الگوریتم به طول می‌انجامد.
	\end{itemize}
	
}

\پاراگراف{سؤال ۲}
پیچیدگی تابع $T(n)$ را در هر قسمت به دست آورید.
\begin{itemize}
	\item $T(n) = 3T(n/3) + \Theta(n(logn)^2)$ [۱۴ نمره]
	\item $T(n) = 5T(n-1) + \Theta(n^2)$ [۱۴ نمره]
	\item $T(n) = T(n - \sqrt{n}) + O(n)$ [۱۴ نمره]
\end{itemize}

\پاسخ{\vsp\vsp}{
	
	\begin{itemize}
		\item 
		درخت بازگشتی 
		$log_3 n = \Theta(logn)$
		تا لایه دارد و هر لایه در مجموع
		$\Theta(n(logn)^2)$
		زمان می‌برد لذا 
		$$T(n) = \Theta(n(logn)^3)$$
		
		\item 
		اگر رابطه را بسط دهیم به عبارت زیر می‌رسیم:
		$$T(n) = \sum_{i=0}^n \Theta(5^i (n-i)^2)$$
		همانطور که می‌بینید عبارت نمایی
		$5^n$
		در نهایت غالب است و 
		$T(n) = \Theta(5^n)$
		را داریم.
		
		\item
		ثابت می‌کنیم تعداد دفعات
		\lr{recursion}
		برابر 
		$\Theta(\sqrt{n})$
		است. داریم که
		$\sqrt{n - \sqrt{n}} \ge \sqrt{n} - 1$
		بنابراین $n$ در مرحله $i$ ام حداقل به اندازه
		$\sqrt{n} - i$
		کم می‌شود لذا در کل در حداکثر
		$\sqrt{n} + 3$
		مرحله 
		\lr{recursion}
		تمام می‌شود. در هر مرحله 
		\lr{recursion}
		ضریبی از 
		$\Theta(n)$
		زمان مصرف شده‌است لذا:
		$$T(n) = \Theta(n \sqrt{n})$$
	\end{itemize}
	
}


\vspace{1pt}
\begin{flushleft}
موفق باشید
\end{flushleft}
\end{document}